\section{Closed-loop control: what an observer must preserve}
\label{sec:design}

Once the controlled model, initial states, allowed intervention rule,
controller, and loss are fixed, only observer errors that affect reachable
states can change the controlled trajectory. An observer turns a model
measurement into information that can guide an intervention. The distinction
between the estimate and the intervention is important. An observer may
correctly report how far the model is from a target while pointing the
controller in an unhelpful direction; conversely, a useful direction can be
paired with a poor estimate.

For model state $h$, let the observer estimate a quantity $\widehat z$ and let
$d$ denote the proposed intervention direction:
\[
E(h)=\widehat z,
\qquad
D(h)=d.
\]
The controller uses the estimate to choose an action magnitude,
\[
u=C(y^\star,\widehat z).
\]
Within this fixed system, testing $E$ keeps $D$ fixed, testing $D$ keeps $E$
fixed, and changing both identifies only the combined system.

\Needspace{6\baselineskip}
An observer need not be correct on every possible model state. It needs to be
correct on the states that the allowed interventions can reach. To state this
concretely, suppose an activation has three coordinates but the controller can
move only the first two. First calibrate the observer to agree with the target
at the starting state. It then needs to predict changes along the first two
coordinates; the third stays fixed. Without that initial calibration, a wrong
coefficient on the unchanged coordinate can still bias every action.

To state the same idea precisely, consider an affine---linear plus an
offset---target and observer:
\[
z(h)=b_T+\beta_T^\top h,
\qquad
\widehat z(h)=b_E+\beta_E^\top h.
\]
Let $h_0$ be the starting state. Let the columns of $B$ span the allowed edit
directions. From $h_0$, the controller can reach only
\[
\mathcal H_R=h_0+\mathcal U,
\qquad
\mathcal U=\operatorname{span}(B).
\]
The observer only has to match the target on this reachable set. With one fixed
direction $d$, this requires both a correct starting estimate and a correct
response to movement along $d$.

\begin{proposition}[Reachable-subspace adequacy]
\label{prop:reachable}
Define
\[
\delta_0=z(h_0)-\widehat z(h_0),
\qquad
\delta_\beta=\beta_T-\beta_E.
\]
Then $z(h)=\widehat z(h)$ for every $h\in\mathcal H_R$ if and only if
\[
\delta_0=0
\qquad\text{and}\qquad
P_{\mathcal U}\delta_\beta=0,
\]
where $P_{\mathcal U}$ is the orthogonal projector onto $\mathcal U$.
\end{proposition}

\begin{proof}
Every reachable state has the form $h=h_0+Bv$. The readout difference is
$\delta_0+\delta_\beta^\top Bv$. It is zero for every $v$ exactly when
$\delta_0=0$ and $B^\top\delta_\beta=0$, which is equivalent to
$P_{\mathcal U}\delta_\beta=0$.
\end{proof}

With one fixed direction $d$, Proposition~\ref{prop:reachable} requires
$\delta_0=0$ and $(\beta_T-\beta_E)^\top d=0$.
The algebra is standard affine feedback analysis. Its role here is to turn
these two requirements into a concrete test of an observer, not to introduce
a new general control theorem.

\subsection{Why a convergent observer can still miss the target}

The proposition tells us where an observer must agree with the target. The
next question is what happens when its remaining error enters a feedback loop.
Consider an unclipped controller that repeatedly edits along one direction:
\[
h_{t+1}=h_t+K\widehat e_t d,
\qquad
\widehat e_t=y^\star-\widehat z(h_t).
\]
Here $y^\star$ is the desired value, $K$ sets the edit size, and
$\widehat e_t$ is the error reported by the observer. Three quantities govern
the loop:
\[
g_E=\beta_E^\top d,
\qquad
g_T=\beta_T^\top d,
\qquad
p=1-Kg_E.
\]
The response $g_E$ is what the observer expects along the edit direction, while
$g_T$ is the target's actual response. The convergence factor $p$---called the
closed-loop pole in control theory---says how quickly the error reported by the
observer shrinks. After $T$ steps, the total action is
\[
A_T=
\begin{cases}
\widehat e_0(1-p^T)/g_E, & g_E\neq0,\\
TK\widehat e_0, & g_E=0,
\end{cases}
\]
and the two errors are exactly
\[
\widehat e_T=p^T\widehat e_0,
\qquad
e_T=e_0-g_T A_T,
\]
where $e_t=y^\star-z(h_t)$. For $g_E\neq0$, define
\[
m_0=z(h_0)-\widehat z(h_0),
\qquad
\rho=\frac{g_T-g_E}{g_E}.
\]
Substitution gives the finite-horizon certificate
\begin{equation}
\label{eq:certificate}
e_T=p^T\widehat e_0-m_0-\rho(1-p^T)\widehat e_0.
\end{equation}

\Needspace{8\baselineskip}
The equation separates three sources of error:

\begin{itemize}
  \item $p^T\widehat e_0$ is the error the observer still reports after $T$
  steps;
  \item $m_0$ is a mismatch at the starting state; and
  \item $\rho(1-p^T)\widehat e_0$ is the mismatch between the response the
  observer expects and the response the target actually has.
\end{itemize}

The observer's reported error converges when $|p|<1$, or $0<Kg_E<2$ for
positive $K$, but this controls only the first term. For example, if $p^T$ is
nearly zero and $m_0=0.2$, the observer reports success while the true target
remains off by about $0.2$, even before response mismatch is counted. A
convergent observer is not necessarily tracking the true target.
